\documentclass[10pt,a4paper]{amsart}
\usepackage[margin=19mm]{geometry}
\usepackage{amsmath,amssymb,mathtools}
\usepackage[scr=rsfs,cal=euler]{mathalfa}
\usepackage{xfrac}
\usepackage[unicode,hidelinks,bookmarks=false]{hyperref}
\newcommand{\ie}{i.e.\ }
\newcommand{\cf}{cf.\ }
\newcommand{\resp}{respectively\ }
\newcommand{\Subsection}{\subsection}
\providecommand{\nobreakdash}{\nobreak\hbox{-}}
\setlength{\emergencystretch}{2em}
\multlinegap0pt
\usepackage{bm}
\usepackage{bbm}
\usepackage{dsfont}
\usepackage{xspace}
\usepackage{cancel}
\usepackage{mathtools}
\usepackage{amsthm}
\makeatletter
\@ifundefined{thm}{\newtheorem{thm}{Theorem}}{}
\@ifundefined{defi}{\newtheorem{defi}[thm]{Definition}}{}
\@ifundefined{lem}{\newtheorem{lem}[thm]{Lemma}}{}
\makeatother
\newcommand{\Q}{\mathbb{Q}}
\newcommand{\ch}{\mathrm{char}}
\renewcommand{\geq}{\geqslant}
\title[Linear relations among radicals]{Linear relations among radicals}
\author{Antonella Perucca}
\date{Source: arXiv:2511.02498v1 (CC BY 4.0)}
\begin{document}
\maketitle

\noindent\emph{Source and licence.} Linear relations among radicals. Version arXiv:2511.02498v1 (CC BY 4.0), distributed under the Creative Commons Attribution 4.0 International licence, as recorded in the arXiv licence metadata for this exact version and shown on the abstract page https://arxiv.org/abs/2511.02498v1. Full rights notice: licenses/arxiv-2511.02498-CC-BY-4.0.txt.\par\smallskip

\noindent\textit{Editorial scope.} Counted unit: the Defi whose source label is \texttt{Delta} in the article (source0.tex lines 407--413, source label \texttt{Delta}), delivered together with the article's own complete original proof of it (source0.tex lines 415--430, one \texttt{proof} environment of 16 lines). No mathematics is written, repaired or re-derived: statement and proof are the article's own text line for line. Required context retained uncounted: mainthm (lines 101--107); Kneserthm (lines 166--170); Ryb2 (lines 282--287); better2 (lines 290--302); lembetter2 (lines 316--326); Schinzel-abelian (lines 387--390). Every definition, display and label of those lines is retained unchanged. Omitted from this package and not redistributed: 1--100, 108--165, 171--281, 288--289, 303--315, 327--386, 391--406, 414--414, 431--528. Nothing inside a retained excerpt is omitted or altered: every retained body is delivered exactly as the article's lines are written, so no omission receipt of any kind accompanies this package. Citations: the retained excerpts carry no citation command at all, so no bibliography entry of the article is transcribed and no reference is redistributed. Reference adaptation: none was needed and none was made; every reference inside the delivered unit resolves inside it, so no unresolved reference or citation is produced. Printed slips kept literal: no printed word, symbol, hypothesis or number of the counted statement or of its proof was altered. Layout and class: the article's own class is \texttt{amsart} with packages including amssymb, amsfonts, amsmath, amsthm, latexsym, babel, version, graphicx, times, titletoc, dirtytalk, inputenc, xcolor, url; the wrapper is the lane's pinned \texttt{amsart} editorial wrapper at 10pt on A4, which restates the theorem-like environments the retained text needs and adds the article's own macro definitions that the retained text needs (\texttt{\textbackslash Q}, \texttt{\textbackslash ch}, \texttt{\textbackslash geq}), with extra wrapper packages (bm, bbm, dsfont, xspace, cancel, mathtools). The delivered numbering is the unit's own. Assets: no retained span contains a figure environment, a graphics-inclusion command, a PDF-page inclusion, a table or a TikZ picture; no image file is copied into this package, the package declares no asset dependency and no image file is redistributed with it. Licence: version arXiv:2511.02498v1 of this article is distributed under the Creative Commons Attribution 4.0 International licence, as recorded in the arXiv licence metadata for this exact version and shown on the abstract page https://arxiv.org/abs/2511.02498v1. A full rights notice is delivered at licenses/arxiv-2511.02498-CC-BY-4.0.txt. Characters: every retained excerpt is pure ASCII, so the delivered engine, XeLaTeX with its default Latin Modern text font, reports no missing character.
\begin{thm}\label{mainthm}
If $n$ is odd, we have
$$\frac{[K(G):K]}{|GK^\times:K^\times|}=\frac{[K(\zeta_{z}):K]}{|\mu_{n}(G K^\times)\cap K(\zeta_{z})^\times: \mu_{n}(K^\times)|}\,.$$
If $n$ is even, writing $n=2^f n'$ where $f$ is a positive integer and $n'$ is an odd integer, the ratio $\frac{[K(G):K]}{|GK^\times:K^\times|}$ equals 
$$\frac{[K(\zeta_{z}):K] \cdot 2^{-\Delta}}{|\mu_{n'}(G K^\times)\cap K(\zeta_{z})^\times: \mu_{n'}( K^\times)|\cdot |\mu_{2^{f+1}}(G K^\times)(G K^\times\cap\sqrt{K^\times})\cap K(\zeta_z)^\times: K^\times|}$$
where $\Delta$ is the non-negative integer from Definition \ref{Delta}.
\end{thm}

\begin{thm}[Kneser's Theorem]\label{Kneserthm}
We have 
$$[K(G):K]=|GK^\times:K^\times|$$
if the following two conditions hold: for every odd prime $p$ we have $\zeta_p\in K^\times$ or $\zeta_p\notin GK^\times$; we have $\zeta_4\in K^\times$ or $1\pm \zeta_4\notin G K^\times$. 
\end{thm}

\begin{thm}[Rybowicz]\label{Ryb2}
We have 
$$\frac{[K(G):K]}{|G^nK^{\times n}:K^{\times n}|}=\delta \cdot [K(\mu_n(GK^\times)):K]$$
where $\delta\in \{1,\frac{1}{2}\}$.
We have $\delta=\frac{1}{2}$ if and only if $a\in G^nK^{\times n}$ and, in the case $w<f$, additionally $1+\zeta_{2^w}\in GK^\times$ and $\ch(K)=0$ and $K\cap \Q(\zeta_{2^\infty})$ is totally real.
\end{thm}

\begin{thm}\label{better2}
Let $\delta$ be as in Theorem \ref{Ryb2} and set
$$H=\begin{cases} \mu_{2n}(GK^\times) & \text{if $w>f$}\\
\mu_{n}(GK^\times) & \text{if $w=f$ and $-\xi_{2^{w+1}}^n\notin G^nK^{\times n}$}\\
\langle 1+\zeta_{2^w}, \zeta_n\rangle & \text{if $w=f$ and $-\xi_{2^{w+1}}^n\in G^nK^{\times n}$}\\
\mu_{n}(GK^\times) & \text{if $w<f$ and $\delta=1$}\\
<1+\zeta_{2^w}>\mu_{n}(GK^\times)& \text{if $w<f$ and $\delta=1/2$}\,.\\
\end{cases}$$
Then $H$ is a subgroup of $GK^\times$ and we have
$$\frac{[K(G):K]}
{|GK^{\times}:K^{\times}|}= \frac{[K(H):K]}
{|HK^{\times}:K^{\times}|}\,.$$
\end{thm}

\begin{lem}\label{lembetter2}
With the notation of Theorem \ref{better2}, we let $m$ be the largest positive integer such that $\zeta_{2^m}\in H$. If $m=1$ then we have 
$\frac{[K(H):K]}{|HK^\times:K^\times|}=1$, while if $m\geq 2$ then we have 
$$\frac{[K(H):K]}{|HK^\times:K^\times|}=\begin{cases} 
2^{2-m} & \text{if $w>f$ or if $w=f$ and $-\xi_{2^{w+1}}^n\notin G^nK^{\times n}$}\\
2^{1-f} & \text{if $w=f$ and $-\xi_{2^{w+1}}^n\in G^nK^{\times n}$}\\
2^{2-\min(w',m)} & \text{if $w<f$, $\delta=1$}\\
2^{1-\min(w,m)} & \text{if $w<f$, $\delta=1/2$}\end{cases}$$
where, if $w$ is finite, $w'$ is the largest positive integer such that $\zeta_{2^{w'}}\in K(\zeta_4)^\times$.
In the last case we have $m=\max(w, \overline{m})$ where $\overline{m}$ is the largest integer such that $\zeta_{2^{\overline{m}}}\in \mu_n(GK^\times)$.
\end{lem}

\begin{thm}[Schinzel's theorem on abelian radical extensions]\label{Schinzel-abelian}
Let $n\geq 1$ be not divisible by $\ch(K)$. 
If $a\in K^\times$, the extension $K(\zeta_{n}, \sqrt[n]{a})/K$ is abelian if and only if $a^{m}=b^{n}$ holds for some $b\in K^{\times}$ and for some $m\mid n$ such that $\zeta_{m}\in K$.
\end{thm}

\begin{defi}\label{Delta}
We set $\Delta=0$ if $\zeta_4\in K^\times$ or $4\nmid n$. In the remaining case, we let 
$H'$ be the group $H$ from Theorem \ref{better2} and Lemma \ref{lembetter2} for $G^{n'}$ over $K(\zeta_z)$ and set 
\begin{equation}\label{magicDelta}
2^{-\Delta}:=\frac{[K(\zeta_z, H'):K(\zeta_z)]}{|H'K(\zeta_z)^\times:K(\zeta_z)^\times|}\,.
\end{equation}
\end{defi}

\begin{proof}[Proof of Theorem \ref{mainthm} if $n$ is even]
Call $G_{P}=\prod_{p\mid n, p\neq 2} G_p$. Remarking that $\zeta_z\in G_P$, we can write     $$\frac{[K(G):K]}{|GK^\times:K^\times|}=[K(\zeta_z):K]\cdot 
\frac{[K(G_{P}):K(\zeta_z)]}{|G_{P}K^\times:K^\times|} \cdot
    \frac{[K(\zeta_z,G_2):K(\zeta_z)]}{|G_2 K^\times:K^\times|}
    $$
By the odd case of Theorem \ref{mainthm} we have
$$\frac{[K(G_{P}):K]}{|G_{P}K^\times:K^\times|}=\frac{[K(\zeta_z):K]}{|\mu_{n/n_2}(G_P K^\times)\cap K(\zeta_{z})^\times: \mu_{n/n_2}( K^\times)|}\,.$$
Since $\mu_{n/n_2}(G K^\times)=\mu_{n/n_2}(G_P K^\times)$ and $\mu_{2n_2}(G K^\times)=\mu_{2n_2}(G_2 K^\times)$ and $GK^\times \cap \sqrt{K^\times}=G_2K^\times \cap \sqrt{K^\times}$ we are left to prove that
$$\frac{[K(\zeta_z,G_2):K(\zeta_z)]}{|G_2 K^\times:K^\times|}=\frac{2^{-\Delta}}{|\mu_{2n_2}(G_2 K^\times)(G_2K^\times\cap\sqrt{K^\times})\cap K(\zeta_z)^\times: K^\times|}\,.$$
As $K(\zeta_z)/K$ is abelian, by Theorem \ref{Schinzel-abelian}
we have   
$$G_2 K^\times\cap K(\zeta_z)^\times=\mu_{2n_2}(G_2 K^\times)(G_2K^\times\cap\sqrt{K^\times})\cap K(\zeta_z)^\times$$
so it suffices to show that 
$$\frac{[K(\zeta_z,G_2):K(\zeta_z)]}{|G_2 K(\zeta_z)^\times:K(\zeta_z)^\times|}=2^{-\Delta}\,,$$
which is a consequence of Theorem \ref{better2} and \eqref{magicDelta} (or of Theorem \ref{Kneserthm} if $\zeta_4\in K^\times$ or $4\nmid n$).
\end{proof}

\end{document}
